\documentclass[../main.tex]{subfiles}

\begin{document}

\section{Eigenvalues and Eigenvectors}

As always, let's start with definitions.

\begin{definition}{}{}
For a linear transformation $T: V \to V$ with a nonzero vector
$v \in V$ and some scalar $\lambda$,if $Tv = \lambda v$, then
$v$ is an \textbf{eigenvector} of $T$ and $\lambda$ is an
\textbf{eigenvalue} of $T$.
\end{definition}

From here, it is important to note that while $v$ must be a
nonzero vector, $\lambda$ could equal zero. Let's take a look
at an example.

\begin{problem}{}{}
Let $T: \mathbb{R}^2 \to \mathbb{R}^2$ be a linear transformation
defined as $T(a, b) = (a, - b)$. Find an eigenvector and eigenvalue
of $T$.
\end{problem}
\begin{solution}
Consider the following equation.
\[
T(0, 1) = (0, -1) = -1 \cdot (0, 1)
\]
Therefore, $(0, 1)$ is an eigenvector of $T$ and $-1$ is the
corresponding eigenvalue of $T$.
\end{solution}

Below is another example for a linear transformation $T: \mathbb{R}^2
\to \mathbb{R}^2$.

\begin{problem}{}{}
Find, if any, eigenvectors and eigenvalues of $T: \mathbb{R}^2 \to
\mathbb{R}^2$ defined as $T(a, b) = (b, -a)$.
\end{problem}
\begin{solution}
For some scalar $\lambda$, consider the following equation.
\[
T(a, b) = (b, -a) = \lambda (a, b) = (\lambda a, \lambda b)
\]
In other words $b = \lambda a$ and $-a = \lambda b$. Substituting,
$-a = \lambda^2 a$ and $(\lambda^2 + 1) a = 0$. Because $\lambda \in
\mathbb{R}$, $a = 0$ and $b = 0$. However, because eigenvectors
are nonzero vectors by definition, $T$ has not real eigenvalues
and eigenvectors.
\end{solution}

Now eigenvectors and eigenvalues are important because they make
our lives easier for investigating matrix representation of a
linear transformation. Consider the following theorem.

\begin{theorem}{}{Eigenvectors and Basis}
For a linear transformation $T: V \to V$ for a finite dimensional
vector space $V$, there exists a basis whose matrix representation
is a diagonal matrix if and only if there exists a basis composed
of only eigenvectors of $T$.
\end{theorem}
\begin{proof}
First, let $B = \{ v_1, \ldots, v_n \}$ be a basis for $V$ and let
the matrix representation of $T$ be a diagonal matrix $A$ whose
diagonal entries are $\lambda_1, \ldots, \lambda_n$. Thus
\[
(Tv_j)_B =
\begin{bmatrix}
    0 \\
    \vdots \\
    \lambda_j \\
    \vdots \\
    0
\end{bmatrix}
\]
hold by construction and $Tv_j = \lambda_j v_j$. In other words,
$v_j$ and $\lambda_j$ are an eigenvector and an eigenvalue of $T$ and
there exists a basis composed of only eigenvectors of $T$.

Continuing, let $B = \{ v_1, \ldots, v_n \}$ be a basis for $V$ whose
elements are all eigenvectors of $T$. Therefore, $Tv_j = \lambda v_j$
and the following equation holds.
\[
(Tv_j)_B =
\begin{bmatrix}
    0 \\
    \vdots \\
    \lambda_j \\
    \vdots \\
    0
\end{bmatrix}
\]
Therefore, there exists a basis whose matrix representation is a
diagonal matrix.
\end{proof}

With this theorem and connecting vector spaces with matrix again,
we can redefine eigenvectors and eigenvalues with different
interpretations.

\begin{definition}{}{}
For a square matrix $A_{n \times n}$ and a nonzero $n$-tuple $x \in
\mathbb{R}^n$, if there exists a scalar $\lambda$ such that
$Ax = \lambda x$, then $x$ is an \textbf{eigenvector} and $\lambda$
is an \textbf{eigenvalue} of $A$.
\end{definition}

When you think about this definition, we can see that both definitions
imply the same thing as the $A$ is the matrix representation of a
linear transformation and $x$ is a vector in the corresponding vector
space. Let's take a look at a quick example.

\begin{problem}{}{}
Find an eigenvector and eigenvalue of matrix $A$ defined as the
following.
\[
A =
\begin{bmatrix}
    1 & 0 \\
    0 & -1
\end{bmatrix}
\]
\end{problem}
\begin{solution}
Consider the following equation.
\[
\begin{bmatrix}
    1 & 0 \\
    0 & -1
\end{bmatrix}
\begin{bmatrix}
    0 \\
    1
\end{bmatrix}
=
\begin{bmatrix}
    0 \\
    -1
\end{bmatrix}
=
-1 \cdot
\begin{bmatrix}
    0 \\
    1
\end{bmatrix}
\]
Thus, $-1$ is an eigenvalue and $[0 1]$ is an eigenvector.
\end{solution}

If this example looks familiar, this is actually the same question
as the first example from this section. Continuing, below are
two interesting theorems on eigenvectors and eigenvalues.

\begin{theorem}{}{Matrix Representation and Eigenvectors}
For a linear transformation $T: V \to V$, $v \in V$ is an
eigenvector with the corresponding eigenvalue $\lambda$ if and only if
$(v)_B$ is an eigenvector with the corresponding eigenvalue $\lambda$
for $A_B^B$, the matrix representation of $T$ for any basis $B$.
\end{theorem}
\begin{proof}
Notice that $Tv = \lambda v$ if and only if $A_B^B (v)_B = (Tv)_B =
(\lambda v)_B = \lambda (v)_B$. Thus, the theorem holds.
\end{proof}

This theorem will come in handy as we can either find all eigenvectors
and corresponding eigenvalues of the transformation or its matrix
representation. Continuing, below is another theorem.

\begin{theorem}{}{Eigenvalues and Subspace}
For a linear transformation $T: V \to V$ and eigenvalue $\lambda$,
$S_\lambda = \{ v \in V | Tv = \lambda v \}$ is a subspace of $V$.
\end{theorem}
\begin{proof}
First, notice that because $S_\lambda \subset V$, it suffices to
show that the set is closed under vector addition and scalar
multiplication. Consider the following equations for $u, v \in
S_\lambda$ and some scalars $\alpha$ and $\beta$.
\begin{align*}
    T(\alpha u + \beta v)
    &= \alpha Tu + \beta Tv
    = \alpha \lambda u + \beta \lambda v \\
    &= \lambda(\alpha u + \beta v)
\end{align*}
In other words, $\alpha u + \beta v \in S_\lambda$ and the theorem
holds.
\end{proof}

As you may have guessed, we can establish a new definition from
the theorem above.

\begin{definition}{}{}
For a transformation $T: V \to V$ and its eigenvalue $\lambda$,
$S_\lambda = \{ v \in V | Tv = \lambda v \}$ defined as the
\textbf{eigenspace} of $T$ for the eigenvalue $\lambda$.
\end{definition}

Just as we redefined eigenvectors and eigenvalues with matrices,
we could do the same for eigenspace. Now that we discussed what
eigenvectors and eigenvalues are, let's discuss how to find one.

\begin{theorem}{}{Eigenvalue and Determinants}
For a square matrix $A_{n \times n}$, $\lambda$ is its eigenvalue
if and only if $\det(A - \lambda I_n) = 0$.
\end{theorem}
\begin{proof}
First notice that by definition, $\lambda$ is an eigenvalue of $A$
if and only if $Ax = \lambda x$ holds for some nonzero vector $x$.
Continuing, $Ax = \lambda x$ holds if and only if $Ax = \lambda I_n x$
and $(A - \lambda I_n)x = 0$. In other words, the columns of the
matrix $A - \lambda I_n$ are linearly dependent and invertible.
Therefore by Theorem \ref{theo:Invertibility and the Determinant},
the theorem holds.
\end{proof}

Being one of the most important theorems in linear algebra,
we can establish another definition from the theorem.

\begin{definition}{}{}
The \textbf{characteristic equation} of a square matrix $A$ is defined
as $\det(A - \lambda I_n) = 0$ for scalar variable $\lambda$.
\end{definition}

An important thing to note from this definition is that
$\det(A - \lambda I_n)$ is a polynomial in $\lambda$ with degree $n$.
Moreover, the real roots of the polynomial can be interpreted as
eigenvalues of $A$. Let's take a look at an example of how
the characteristic equation can be used to find all eigenvalues of
a matrix.

\begin{problem}{}{}
Find all eigenvalues and the corresponding eigenspaces of matrix $A =
\begin{bmatrix}
    1 & 0 \\
    0 & -1
\end{bmatrix}$.
\end{problem}
\begin{solution}
To find all eigenvalues of $A$, it suffices to find all solutions for
the characteristic equation of $A$ by Theorem \ref{theo:Eigenvalue and
Determinants}.
\[
A - \lambda I_2 =
\begin{bmatrix}
    1 - \lambda & 0 \\
    0           & -1 - \lambda
\end{bmatrix}
\]
Therefore, $\det(A - \lambda I_2) = (1 - \lambda)(-1 - \lambda)$
and $\lambda = 1$ or $\lambda = -1$ are the solutions to the
characteristic equation.

Continuing, the corresponding eigenspaces could be found. For
$\lambda = 1$, the eigenspace is defined as $S_1 = \{ x | Ax = x \}$.
Solving for the equation $(A - I)x = 0$, the following equations hold.
\begin{align*}
    \begin{bmatrix}
        0 & 0 \\
        0 & -2
    \end{bmatrix}
    \begin{bmatrix}
        x_1 \\
        x_2
    \end{bmatrix}
    &= 0 \\
    -2x_2 &= 0
\end{align*}
Therefore, the eigenspace for the eigenvalue $1$ can be represented
as the following
\[
S_1 =
\Span\left\{
    \begin{bmatrix}
        1 \\
        0
    \end{bmatrix}
\right\}
\]
Similarly, $S_{-1} = \{ x | Ax = -x \}$ and $(A + I)x = 0$ could
be solved.
\begin{align*}
    \begin{bmatrix}
        2 & 0 \\
        0 & 0
    \end{bmatrix}
    \begin{bmatrix}
        x_1 \\
        x_2
    \end{bmatrix}
    &= 0 \\
    2x_1 &= 0
\end{align*}
Therefore,
\[
S_{-1} =
\Span\left\{
    \begin{bmatrix}
        0 \\
        1
    \end{bmatrix}
\right\}
\]
and the corresponding eigenspaces for all eigenvalues of $A$ are
found.
\end{solution}

We got a nice matrix here, but it's worth noting that $\lambda$ need
not be real. As demonstrated from the simple example above,
solving the characteristic equation will lead to all eigenvalues
for a matrix. I will leave more complex and interesting examples
in the practice problems below. In the next part of the section,
we will discuss the properties of eigenvalues and eigenvectors.

\subsection{Properties of Eigenvalues}

Eigenvalues and eigenvectors have interesting properties and
relationship with vector spaces. Consider the following theorem.

\begin{theorem}{}{Similar Matrices and Eigenvalues}
For two square matrices $A$ and $B$ such that $A \cong B$,
the characteristic equation and eigenvalues for each matrix are equal.
\end{theorem}
\begin{proof}
By definition of similar matrices, there exists some invertible
matrix $P$ such that $B = P^{-1} A P$. Therefore, the following
equations hold.
\[
P^{-1} (A - \lambda I) P
= P^{-1} A P - P^{-1} \lambda I P
= B - \lambda I
\]
Therefore, $\det(B - \lambda I) = \det(P^{-1} (A - \lambda I) P)
= \det(A - \lambda I)$ and the characteristic equations are the
same. In other words, the eigenvalues for $A$ and $B$ are equal
and the theorem holds.
\end{proof}

Below is another theorem on the eigenvalues of triangular matrices.

\begin{theorem}{}{Triangular Matrices and Eigenvalues}
The eigenvalues of triangular (upper or lower) matrices are the
diagonal elements.
\end{theorem}
\begin{proof}
By definition, notice that for a triangular matrix $A$,
$A - \lambda I$ is also triangular. Therefore by Theorem
\ref{theo:The Determinants of Triangular Matrices},
\[
\det(A - \lambda I)
= \prod_{k=1}^n (\lambda_k - \lambda)
\]
and $\det(A - \lambda I_n) = 0$ if and only if $\lambda = \lambda_k$.
Because $\lambda_k$ are the diagonal elements of $A$, the theorem
holds.
\end{proof}

Continuing, below is another theorem on the relationship between
eigenvalues and invertibility.

\begin{theorem}{}{Invertibility and Eigenvalues}
A square matrix $A$ contains a zero eigenvalue if and only if it is
not invertible.
\end{theorem}
\begin{proof}
A square matrix $A$ contains a zero eigenvalue if and only if
$\det(A - 0 I) = 0$ and $\det(A) = 0$. Because $A$ is invertible
if and only if $\det(A) \neq 0$ by Theorem \ref{theo:Invertibility
and the Determinant}, the theorem holds.
\end{proof}

If you think this theorem is nice, below is another theorem on
invertibility and eigenvalues.

\begin{theorem}{}{Eigenvalues and Eigenvectors of an Invertible Matrix}
For an invertible matrix $A$ with eigenvector $x$ and its
corresponding eigenvalue $\lambda$, $x$ is also an eigenvalue of
$A^{-1}$ with corresponding eigenvalue $\frac{1}{\lambda}$.
\end{theorem}
\begin{proof}
By definition, $Ax = \lambda x$ and $x = A^{-1} \lambda x$.
Continuing, $A^{-1} x = \frac{1}{\lambda} \cdot x$. Therefore,
$x$ is also an eigenvector of $A^{-1}$ with the corresponding
eigenvalue $\frac{1}{\lambda}$.
\end{proof}

Continuing, below is a theorem on the relationship between
eigenvalues and linear independence.

\begin{theorem}{}{Eigenvalues and Linear Independence}
If $x_1, \ldots, x_k$ are eigenvectors of a square matrix $A$
and the corresponding eigenvalues $\lambda_1, \ldots, \lambda_k$
are all distinct, then $\{ x_1, \ldots, x_k \}$ is linearly
independent.
\end{theorem}
\begin{proof}
First, notice that $\{ x_1 \}$ is linearly independent. Therefore,
it suffices to show that if the theorem holds for $k - 1 > 0$, then
the statement holds for $k$. Assume $\{ x_1, \ldots, x_{k-1} \}$
is a linearly independent set of eigenvectors with distinct
corresponding eigenvalues $\lambda_1, \ldots, \lambda_{k-1}$.
Consider the following equations for some scalars $c_i$.
\begin{align*}
    \sum_{i=1}^k c_i x_i &= 0 \\
    \sum_{i=1}^k c_i A x_i &= 0 \\
    \sum_{i=1}^k c_i \lambda_i x_i &= 0 \\
    \sum_{i=1}^{k-1} c_i \lambda_k x_i &= 0
\end{align*}
Subtracting the last two equations, $\sum_{i=1}^{k-1} c_i
(\lambda_i - \lambda_k) x_i = 0$ is obtained. Because the set
$\{ x_1, \ldots, x_{k-1} \}$ is linearly independent, $c_i
(\lambda_i - \lambda_k) = 0$ for all integer $i \in [1, k)$.
Continuing, because $\lambda_i - \lambda_k \neq 0$ by construction,
$c_i = 0$ for all integer $i \in [1, k)$ and $c_k = 0$ since
$x_k \neq 0$. Thus $\{ x_1, \ldots, x_k \}$ is linearly independent
and the theorem holds by induction.
\end{proof}

Naturally, we can consider the dimension of an eigenspace.

\begin{theorem}{}{The Dimension of an Eigenspace}
For a matrix $A_{n \times n}$ with an eigenvalue $\lambda$,
the following equation holds.
\[
\dim(S_\lambda) = n - \Rank(A - \lambda I)
\]
\end{theorem}
\begin{proof}
By definition, the eigenspace for $\lambda$ is the set of all
vectors $v$ that satisfy $Av = \lambda v$. Because $(A - \lambda) v =
0$, $S_\lambda$ can be interpreted as $\ker(A - \lambda I)$.
Moreover by Theorem \ref{theo:Rank and Nullity of T}, the following
equation holds.
\[
\Rank(A - \lambda I) + \Null(A - \lambda I) = n
\]
Because $\dim(\ker(A - \lambda I)) = \Null(A - \lambda I)$ holds by
definition, $\dim(S_\lambda) = n - \Rank(A - \lambda I)$.
\end{proof}

We could really go on forever with interesting theorems, but
before we conclude this part of the note, I would like to discuss
two more theorems.

\begin{theorem}{}{Properties of Eigenvector and Eigenvalue}
Consider the following statements for a matrix $A$ with an
eigenvector $x$ with its corresponding eigenvalue $\lambda$.
\begin{enumerate}[noitemsep]
\item For an arbitrary scalar $k$, $x$ is also an eigenvector
of $kA$ with corresponding eigenvalue $k \lambda$.
\item For some $n \in \mathbb{N}$, $x$ is also an eigenvector
of $A^n$ with corresponding eigenvalue $\lambda^n$.
\end{enumerate}
\end{theorem}

Let's start by proving the first statement.

\begin{proof}
Consider the following equation.
\[
(kA)x = k(Ax) = k \lambda x = (k \lambda) x
\]
Therefore, $x$ is an eigenvector of $kA$ and its corresponding
eigenvalue is $k \lambda$.
\end{proof}

Below is the proof for the second statement.

\begin{proof}
Assume that if the statement holds for $n - 1$, then the statement
holds for $n > 1$. Consider the following equations.
\[
A^n x
= A \left( A^{n-1} x \right)
= A \left( \lambda^{n-1} x \right)
= \lambda^{n-1} Ax
= \lambda^{n-1} \lambda x
= \lambda^n x
\]
Because $Ax = \lambda x$, the statement holds by induction.
\end{proof}

Before we move on to the final content with the last theorem for this
part, below is a definition for the theorem.

\begin{definition}{}{}
The \textbf{trace} of a matrix $A = [a_{ij}]$, denoted as $\Tr(A)$,
is the sum of the diagonal elements of $A$.
\end{definition}

With the definition, we can continue with the theorem.

\begin{theorem}{}{Properties of Eigenvalues}
For a square matrix $A_{n \times n}$ with eigenvalues $\lambda_1,
\ldots, \lambda_n$, the following equations hold.
\begin{enumerate}[noitemsep]
\item $\Tr(A) = \sum_{k=1}^n \lambda_k$
\item $\det(A) = \prod_{k=1}^n \lambda_k$
\end{enumerate}
\end{theorem}

Let's start by proving the first equation.

\begin{proof}
By the Fundamental Theorem of Algebra, the following equation holds.
\[
\det(A - \lambda I)
= \prod_{k=1}^n (\lambda_k - \lambda)
= 0
\]
Therefore, all eigenvalues are the diagonal elements and by the
definition of trace of $A$, the sum of all eigenvalues are the
trace of $A$.
\end{proof}

Continuing, below is the proof for the second equation.

\begin{proof}
Substituting $\lambda = 0$ into the equation
\[
\det(A - \lambda I)
= \prod_{k=1}^n (\lambda_k - \lambda),
\]
$\det(A) = \prod_{k=1}^n \lambda_k$ and the equation holds.
\end{proof}

It is important to note that the second equation counts multiplicity.
In other words, even if $\lambda_p = \lambda_q$, you still count
both. Wow! I can't believe that we are already almost done with the
notes! For the final concept of my notes on linear algebra, we will
discuss diagonalization.

\subsection{Diagonalization}

Diagonalization in linear algebra is a technique that I think
really makes our lives easier in calculations. To get into it,
let's get started with definition and theorems.

\begin{definition}{}{}
A square matrix that is similar to a diagonal matrix is known to be
\textbf{diagonalizable}.
\end{definition}

Now we cannot always use a brute force method to see if two matrices
are similar. Fortunately, the following theorem will help in checking
whether a square matrix is diagonalizable.

\begin{theorem}{}{Diagonalizability and Eigenvectors}
A square matrix $A_{n \times n}$ is diagonalizable if and only if
it retains a set of $n$ eigenvectors that is linearly independent.
\end{theorem}
\begin{proof}
First and foremost, notice that $A$ and a diagonal matrix $D$ of
the same order with entries $\lambda_i$ are similar if and only if
there exists an invertible matrix $P$ such that $P^{-1} A P = D$ or
$AP = PD$. Let $x_j$ and $e_j$ be the $j^\text{th}$ column of $P$ and
$I$ respectively. Then, the following equations hold.
\[
A x_j
= A Pe_j
= PDe_j
= P \lambda_j e_j
= \lambda_j P e_j
= \lambda_j x_j
\]
Thus, $x_j$ and $\lambda_j$ are an eigenvector and the corresponding
eigenvalue of $A$. By definition, because $M$ is invertible,
$x_j$ are all linearly independent and $A$ has $n$ linearly
independent eigenvectors. In other words, if $A$ is diagonalizable,
then $A$ has $n$ linearly independent eigenvectors.

Continuing, the converse can be proven. Let $A$ have $n$ linearly
independent eigenvectors denoted as $x_j$ with corresponding
eigenvalues $\lambda_j$. Construct matrices $P$ and $D$ such that
the $j^\text{th}$ column of $P$ is $x_j$ and the diagonal element
of $D$ is $\lambda_j$. If $e_j$ is the $j^\text{th}$ column of $I$,
then the $j^\text{th}$ column of $AP$ can be written as the following.
\[
APe_j = A x_j = \lambda_j x_j
\]
Similarly, the $j^\text{th}$ column of $PD$ can be written as $PDe_j =
M \lambda_j e_j = \lambda_j Me_j = \lambda_j x_j$ and $AP = PD$ holds.
Moreover, because $P$ is invertible by construction, $P^{-1} A P = D$
and $A$ is diagonalizable.
\end{proof}

Continuing from the theorem, we can establish the following corollary.

\begin{corollary}{}{Corollary of Diagonalizability and Eigenvectors}
A matrix $A_{n \times n}$ with $n$ distinct real eigenvalues is
diagonalizable.
\end{corollary}
\begin{proof}
First and foremost, notice that by Theorem \ref{theo:Eigenvalues and
Linear Independence}, the set of corresponding eigenvectors
$\{ x_1, \ldots, x_n \}$ of $A$ is linearly independent. Thus by
Theorem \ref{theo:Diagonalizability and Eigenvectors}, the corollary
holds.
\end{proof}

This theorem also naturally leads to the question if a square matrix
can be similar to more than one diagonal matrix. The answer to
that question can be found with the theorem below.

\begin{theorem}{}{Multiple Diagonal Matrices}
If a square matrix is similar to both diagonal matrices $D_1$ and
$D_2$, then the diagonal entries of $D_1$ and $D_2$ are the same
with possibly different orderings.
\end{theorem}
\begin{proof}
First and foremost, notice that by Theorem \ref{theo:Similar Matrices
and Eigenvalues}, because $A$ and $D_1$ are similar, they have the
same characteristic equation. Similarly, because $A$ and $D_2$ are
similar, they also have the same characteristic equation and the
characteristic equations of $D_1$ and $D_2$ are equal.
\[
\det(D_1 - \lambda I) = \det(D_2 - \lambda I) = 0
\]
By construction, notice that $D_1 - \lambda I$ and $D_2 - \lambda I$
are diagonal matrices. Thus by Corollary \ref{cor:Determinants of
Diagonal Matrices}, the determinant is the product of all $(\lambda
- \lambda_k$ for integer $k \in [1, n]$ where $\lambda_k$ are the
diagonal entries for $D_1$. Because multiplication is commutative,
the diagonal entries for $D_1$ and $D_2$ are the same.
\end{proof}

Naturally, we can see that if $A \cong D_1$ and $A \cong D_2$,
then $D_1 \cong A$ and $D_1 \cong D_2$. In other words, if $A$ is
similar to multiple diagonal matrices, then the diagonal matrices
are similar to each other.

Now that we discussed when a matrix is diagonalizable, let's discuss
how to diagonalize and its applications.

\begin{definition}{}{}
\textbf{Diagonalization} is the process of writing a diagonalizable
matrix $A$ in the form of $A = P D P^{-1}$.
\end{definition}

Diagonalization is particularly useful as it makes calculations
easier. Before concluding our notes, below is the final theorem!

\begin{theorem}{}{Diagonalizability and Eigenspace}
Let $A_{n \times n}$ be a matrix with real and distinct eigenvalues
$\lambda_1, \ldots, \lambda_k$. The matrix $A$ is diagonalizable if
and only if the following equation holds.
\[
\sum_{i=1}^k \dim(S_{\lambda_i}) = n
\]
\end{theorem}
\begin{proof}
First, the statement if $\sum_{i=1}^k \dim(S_{\lambda_i}) = n$,
then $A$ is diagonalizable can be proven. Let $B_i$ be a basis
for eigenspace $S_{\lambda_i}$. By definition, vectors $v$ in
eigenspace $S_{\lambda_i}$ satisfy $Av = \lambda_i v$. Because
$\lambda_i$ are all distinct, no eigenspace shares the same
eigenvector and thus no basis of two different eigenspaces
shares a vector. In other words, $B$ contains $n$ distinct eigenvectors
of $A$. Let $v_i \in S_{\lambda_i}$ be some linear combination of
vectors in $B_i$. Notice that if $\sum_{i=1}^k v_i = 0$, then the
coefficients must be zero since $\{ v_1, \ldots, v_k \}$ is linearly
independent by Theorem \ref{theo:Eigenvalues and Linear Independence}.
In other words, all coefficients for the linear combination of $B_i$
that forms $v_i$ must be zero for $\sum_{i=1}^k v_i = 0$ to hold.
Thus, $B$ is linearly independent and $A$ is diagonalizable by
Theorem \ref{theo:Diagonalizability and Eigenvectors}.

Continuing, the converse can be proven. Assume that $\sum_{i=1}^k
\dim(S_{\lambda_i}) > n$. Continuing from the previous definition,
$B \subset \mathbb{R}^n$ and the inequality $\sum_{i=1}^k
\dim(S_{\lambda_i}) > n$ leads to contradiction. Moreover, if
$\sum_{i=1}^k \dim(S_{\lambda_i}) < n$ then $A$ has no set that
contains $n$ linearly independent eigenvectors. Thus, $A$ is not
diagonalizable if $\sum_{i=1}^k \dim(S_{\lambda_i}) \neq n$.
\end{proof}

This is it for eigenvalues and eigenvectors! This note was a
particularly long one, but I think it also means that the
relationships between matrices and vector spaces are one of the
most essential understandings for linear algebra. Below are a few
practice problems on the topics that we discussed.

\end{document}


